\section{Signal Processing Pipeline for Passive ELINT}\label{sec:bg:elint_pipeline}

In \gls{ew} and \gls{esm} applications, measured radar energy passes through a staged \emph{signal processing pipeline} before it reaches higher-level \gls{elint} reasoning.
The exact implementation varies with platform, band, and mission software, but the logical chain is broadly stable across textbooks and system descriptions.

At the \emph{front end}, antennas and \gls{rf} conditioning capture emissions over one or more instantaneous bandwidths, often after wideband digitization or channelized sub-band processing.
A \emph{detection} stage then marks pulse-like events against a noise and clutter background, producing a sparse list of candidate pulse times (and sometimes coarse frequency hints).

Given detections, \emph{parameter estimation} attaches per-pulse measurements such as carrier frequency, pulse width, time of arrival, direction of arrival when available, and amplitude or signal-to-noise proxies.
These estimates are noisy, can be missing for weak emitters, and may include false alarms that are not associated with any physical radar.

\emph{Sorting} and \emph{deinterleaving} group pulses into coherent pulse trains by exploiting periodic structure in inter-pulse timing and consistent parameter tracks across time.
The output of this stage is commonly packaged as a stream of \glspl{pdw} (or closely related pulse descriptor records), which form the usual interface to downstream tasks such as emitter identification and mode analysis (\cref{sec:bg:pdw}).

In classical \gls{elint} architectures, specifications therefore treat train formation as an early, largely discrete step: deinterleaving is completed first so that downstream logic receives one dominant hypothesis per physical emission path, or a short ranked list, rather than a time-interleaved mixture \parencite{wiley2006elint}.
\emph{Emitter reports} are assembled only after that gate, summarizing each train with aggregated \gls{pdw} statistics, tentative equipment identities from lookup tables, and temporal extent.
Mode classification is then commonly framed as matching the train against a \emph{mode library} of nominal \gls{rf} bands, \gls{pri}/\gls{prf} families, stagger laws, and scan programs, so that declared modes remain anchored to catalogued waveform templates.
\emph{Positional estimation}---for example \gls{aoa} intersections from direction-finding networks, time-difference-of-arrival fixes along baselines, or fused tracks across collectors---consumes the same deinterleaved associations, which isolates geometry and timing errors from the sorting stage.

Later stages may maintain libraries, threat lists, and operator displays, but for the learning problem in this thesis the critical contract is that the model consumes the deinterleaved \gls{pdw} stream produced here.

% TODO: Add a schematic figure (figures/background/) and cite a standard EW/ESM text for the pipeline stages.
